% ju 28-4-22 
\documentclass[a4paper,fontsize=14pt,DIV=calc,fleqn,parskip=half]{scrartcl}
\usepackage[ngerman]{babel}
\usepackage[utf8]{inputenc}
\usepackage[T1]{fontenc}

\usepackage{lmodern}
\usepackage[osf,sc]{mathpazo} 
\usepackage[scale=.9,semibold]{sourcecodepro}   
%\usepackage[osf]{sourcesanspro}   

%\usepackage[left=2cm,right=2cm,top=1cm,bottom=1cm,includeheadfoot]{geometry}
%\usepackage[left=4cm,right=2cm,top=1cm, bottom=1cm,includeheadfoot]{geometry}
%\usepackage[left=6cm,right=1cm,top=1cm, bottom=1cm,includeheadfoot]{geometry}
\usepackage[landscape=true,left=2cm,right=2cm,top=1cm,bottom=1cm,includeheadfoot]{geometry}%quer
\usepackage{scrlayer-scrpage}
\pagestyle{scrheadings}
\clearpairofpagestyles
\addtokomafont{section}{\color{blau5}}

\usepackage{xcolor}
% eigene Farbe definieren
% Adobe Prozessfarben: CMYK: 100,50,0,35 -> 1,0.5,0,0.35
\definecolor{orange}{cmyk}{0,0.55,0.61,0}   % 0,55,61,0
\definecolor{blau5}{cmyk}{1,0.77,0.1,0.01}  % 100,77,10,
\definecolor{rot5}{cmyk}{0.22,1,1,0.19}     % 22,100,100,19
\definecolor{grau2}{cmyk}{0,0,0,0.1}        % 0,0,0,40
\definecolor{blau}{cmyk}{0.93,0.66,0,0.21}% 
\usepackage{amsmath}
\usepackage{amsfonts}
\usepackage{amssymb}
\usepackage{graphicx}
\usepackage{subcaption}
\usepackage{booktabs}
%\usepackage{blindtext}
\usepackage{listings}
\usepackage{eurosym}%\euro 20,-
\usepackage{url}% Links
\usepackage[inline]{enumitem}
\usepackage{pifont}

%%%%%%%%%%%%%%%%%%%%%%%%%%%%%%%%%%%%%%%%%%%%%%%%%%%%%%%
\newcommand{\name}{Jan Unger}% anpassen!!!!!
\newcommand{\thema}{\LaTeX -Vorlage}% anpassen!!!!!
\newcommand{\quelle}{\name}
%%%%%%%%%%%%%%%%%%%%%%%%%%%%%%%%%%%%%%%%%%%%%%%%%%%%%%%

%\ihead{\textbf{Quelle:} \quelle}%{Kopfzeile innen}
%\ohead{\textbf{Datum:} \today}  %{Kopfzeile außen}
\ifoot{\textbf{Thema:} \thema}  %{Fußzeile  innen}
\ofoot{\pagemark}             %{Fußzeile  außen}

\title{\color{blau5}{\thema}}
\author{\name}
\date{\today}

\begin{document}

\thispagestyle{empty}
\maketitle
\newpage
\setcounter{page}{1}

\section*{Listen}%\label{sec:listen}\index{Listen}

% Checkliste
\begin{itemize}[label=\checkmark] %\itemsep -2pt
	\item Check
\end{itemize}

% Checkliste
\begin{itemize} 
	\item [$\square$] Check
\end{itemize}

% Liste
\begin{itemize} 
	\item Punkt
\end{itemize}
% Liste numeriert
\begin{enumerate} 
	\item Liste numeriert
	\item Punkt
\end{enumerate}

\newpage
\section*{Abbildungen}%\label{sec:Abbildungen}\index{Abbildungen}

Liste

\begin{itemize} 
	\item Punkt
	\item Punkt
	\item Punkt
	\item Punkt
\end{itemize}
		
\begin{figure}[!h]
\centering
\subcaptionbox{Logo 1 \label{logo1}}
{\includegraphics[width=0.25\textwidth]{Logo/Logo1}}
\subcaptionbox{Logo 2 \label{logo2}}
{\includegraphics[width=0.25\textwidth]{Logo/Logo2}}
\caption{Zwei Logos}\label{Logos}
\end{figure}

\newpage
\section*{Abbildungen 2}%\label{sec:Abbildungen2}\index{Abbildungen2}
%Logo in Neg, Grau, Schwarz (\autoref{fig:logoneggrauschwarz}).
%
\begin{figure}[!h]% hier: !hb
	\centering
	\begin{minipage}[b]{0.40\textwidth}
		\includegraphics[width=\textwidth]{Logo/Logo-negativ}%
	\end{minipage}
	\hfill
	\begin{minipage}[b]{0.30\textwidth}
		\includegraphics[width=\textwidth]{Logo/Logo-Grau}%
	\end{minipage}
	\hfill
	\begin{minipage}[b]{0.20\textwidth}
		\includegraphics[width=\textwidth]{Logo/Logo-SW}%
	\end{minipage}
	\caption{Logo in Neg, Grau, Schwarz}\label{fig:logoneggrauschwarz}%% anpassen
\end{figure}

\newpage
\section*{Text und Abbildungen}%\label{sec:TAbbildungen}\index{Text, Abbildungen}
%Logo in Neg, Grau, Schwarz (\autoref{fig:logoneggrauschwarz}).
%
\begin{figure}[!h]% hier: !hb
	\centering
	\begin{minipage}[c]{0.55\textwidth}
		Liste
		\begin{itemize} 
			\item Punkt
			\item Punkt
			\item Punkt
			\item Punkt
		\end{itemize}
	\end{minipage}
	\hfill
	\begin{minipage}[c]{0.35\textwidth}
		\includegraphics[width=\textwidth]{Logo/Logo2}%
	\end{minipage}
	%\caption{Logo in Neg, Grau, Schwarz}\label{fig:logoneggrauschwarz}%% anpassen
\end{figure}

\end{document}
